% =====================================================================
% CAP. 11 - NIVEL 2
% =====================================================================
\blocoaplicacao

\begin{questao}[2]{}
A tensão de um capacitor muda de $\SI{8}{\volt}$ para um valor final de
$\SI{2}{\volt}$ através de $R=\SI{4}{\kilo\ohm}$ e
$C=\SI{25}{\micro\farad}$. Determine $v_C(\SI{0.2}{\second})$.
\resposta{$v_C(\SI{0.2}{\second})\approx\SI{2.81}{\volt}$}
\resolucao{$\tau=RC=\SI{0.1}{\second}$. Logo
\[
v_C(t)=2+(8-2)e^{-t/0{,}1}.
\]
Em $t=\SI{0.2}{\second}$, $v_C=2+6e^{-2}\approx\SI{2.81}{\volt}$.}
\end{questao}

\begin{questao}[2]{}
A corrente de um indutor passa de $\SI{1}{\ampere}$ para um valor final de
$\SI{4}{\ampere}$. Sabendo que $L=\SI{0.6}{\henry}$ e
$R_{\mathrm{eq}}=\SI{30}{\ohm}$, calcule $i_L(\SI{40}{\milli\second})$.
\resposta{$i_L\approx\SI{3.59}{\ampere}$}
\resolucao{$\tau=L/R=0{,}6/30=\SI{20}{\milli\second}$. Assim
$i_L(t)=4+(1-4)e^{-t/\tau}$. Para $t=2\tau$,
$i_L=4-3e^{-2}\approx\SI{3.59}{\ampere}$.}
\end{questao}

\begin{questao}[2]{}
Um circuito RC inicialmente descarregado possui $R=\SI{10}{\kilo\ohm}$ e
$C=\SI{10}{\micro\farad}$. Quanto tempo é necessário para atingir
\SI{90}{\percent} do valor final?
\resposta{$t\approx\SI{230}{\milli\second}$}
\resolucao{$\tau=\SI{0.1}{\second}$. De
$0{,}9=1-e^{-t/\tau}$ resulta $e^{-t/\tau}=0{,}1$ e
$t=\tau\ln 10\approx\SI{0.230}{\second}$.}
\end{questao}

\begin{questao}[2]{}
Um capacitor descarrega de $\SI{12}{\volt}$ para $\SI{3}{\volt}$ com
$\tau=\SI{50}{\milli\second}$. Determine o tempo decorrido.
\resposta{$t\approx\SI{69.3}{\milli\second}$}
\resolucao{$3=12e^{-t/\tau}$, então $e^{-t/\tau}=1/4$ e
$t=\tau\ln4=\SI{50}{\milli\second}\cdot1{,}386\approx\SI{69.3}{\milli\second}$.}
\end{questao}

\begin{questao}[2]{}
A corrente natural de um circuito RL cai de $\SI{2}{\ampere}$ para
$\SI{0.2}{\ampere}$. Se $\tau=\SI{25}{\milli\second}$, determine o intervalo de
tempo.
\resposta{$t\approx\SI{57.6}{\milli\second}$}
\resolucao{$0{,}2=2e^{-t/\tau}$ implica $e^{-t/\tau}=0{,}1$. Logo
$t=\tau\ln10\approx25\cdot2{,}303=\SI{57.6}{\milli\second}$.}
\end{questao}

\begin{questao}[2]{}
No circuito abaixo, a chave conecta a fonte em $t=0$ e o capacitor está
inicialmente descarregado. Determine $v_C(\tau)$ e a constante de tempo.

\circuitoq{%
\begin{circuitikz}[scale=1,transform shape,line width=0.8pt]
\draw (0,0) to[battery1,l_={$\SI{20}{\volt}$}] (0,3)
      to[R,l={$R_1=\SI{6}{\kilo\ohm}$}] (3,3) coordinate(a);
\draw (a) to[R,l={$R_2=\SI{3}{\kilo\ohm}$}] (3,0) -- (0,0);
\draw (a) -- (5,3) to[C,l={$C=\SI{50}{\micro\farad}$}] (5,0) -- (3,0);
\end{circuitikz}}

\resposta{$\tau=\SI{100}{\milli\second}$; $v_C(\tau)\approx\SI{4.21}{\volt}$}
\resolucao{Em regime permanente, a tensão no capacitor corresponde à tensão do divisor:
$V_f=20\cdot3/(6+3)=\SI{6.67}{\volt}$. Suprimindo a fonte,
$R_{\mathrm{eq}}=R_1\parallel R_2=\SI{2}{\kilo\ohm}$, então
$\tau=R_{\mathrm{eq}}C=\SI{0.1}{\second}$. Como $v_C(0)=0$,
$v_C(\tau)=V_f(1-e^{-1})\approx\SI{4.21}{\volt}$.}
\end{questao}

\begin{questao}[2]{}
Após uma comutação, um capacitor possui $v_C(0^+)=\SI{5}{\volt}$ e tende a
$-\SI{10}{\volt}$, com $\tau=\SI{20}{\milli\second}$. Determine
$v_C(\SI{10}{\milli\second})$.
\resposta{$v_C\approx-\SI{0.90}{\volt}$}
\resolucao{$v_C(t)=-10+[5-(-10)]e^{-t/\tau}=-10+15e^{-t/\tau}$.
Para $t/\tau=0{,}5$, obtém-se $v_C\approx-10+15e^{-0{,}5}
\approx-\SI{0.90}{\volt}$.}
\end{questao}

\begin{questao}[2]{}
Um indutor de $\SI{40}{\milli\henry}$ está conectado, após uma
comutação, a uma rede cuja resistência equivalente vista em seus terminais é
$R_{\mathrm{eq}}=\SI{4}{\ohm}$. Sua corrente inicial é $\SI{0.5}{\ampere}$ e o valor final é
$\SI{2}{\ampere}$. Calcule a corrente após $\SI{20}{\milli\second}$.
\resposta{$i_L\approx\SI{1.80}{\ampere}$}
\resolucao{$\tau=L/R=\SI{10}{\milli\second}$ e
$i_L(t)=2+(0{,}5-2)e^{-t/\tau}$. Em $t=2\tau$,
$i_L=2-1{,}5e^{-2}\approx\SI{1.80}{\ampere}$.}
\end{questao}

\begin{questao}[2]{}
Um circuito RLC em série tem $R=\SI{10}{\ohm}$, $L=\SI{50}{\milli\henry}$ e
$C=\SI{200}{\micro\farad}$. Determine $\alpha$, $\omega_0$, $\omega_d$ e o tipo
de resposta.
\resposta{$\alpha=\SI{100}{\per\second}$; $\omega_0\approx\SI{316.2}{\radian\per\second}$;
$\omega_d=\SI{300}{\radian\per\second}$; subamortecida}
\resolucao{$\alpha=R/(2L)=10/0{,}1=100$. Além disso,
$\omega_0=1/\sqrt{0{,}05\cdot\pot{200}{-6}}\approx316{,}2$. Como
$\alpha<\omega_0$, a resposta é subamortecida e
$\omega_d=\sqrt{316{,}2^2-100^2}=\SI{300}{\radian\per\second}$.}
\end{questao}

\begin{questao}[2]{}
Para um circuito RLC em série com $R=\SI{40}{\ohm}$, $L=\SI{50}{\milli\henry}$ e
$C=\SI{200}{\micro\farad}$, classifique a resposta natural.
\resposta{Superamortecida}
\resolucao{$\alpha=40/(2\cdot0{,}05)=\SI{400}{\per\second}$, enquanto
$\omega_0\approx\SI{316.2}{\radian\per\second}$. Como $\alpha>\omega_0$, a
resposta é superamortecida.}
\end{questao}

\begin{questao}[2]{}
Calcule $R_{\mathrm{crit}}$ para um circuito RLC em série com $L=\SI{0.1}{\henry}$ e
$C=\SI{100}{\micro\farad}$.
\resposta{$R_{\mathrm{crit}}\approx\SI{63.25}{\ohm}$}
\resolucao{$R_{\mathrm{crit}}=2\sqrt{L/C}=2\sqrt{0{,}1/\pot{100}{-6}}
=2\sqrt{1000}\approx\SI{63.25}{\ohm}$.}
\end{questao}

\begin{questao}[2]{}
Uma resposta subamortecida de tensão tem a forma
\[
v(t)=e^{-100t}\big(10\cos300t+B\sin300t\big).
\]
Sabendo que $\mathrm{d}v/\mathrm{d}t|_{0^+}=0$, determine $B$.
\resposta{$B\approx\SI{3.33}{\volt}$}
\resolucao{Derivando e avaliando em zero,
$\dot v(0)=-100\cdot10+300B=0$. Portanto
$B=1000/300\approx\SI{3.33}{\volt}$.}
\end{questao}

\begin{questao}[2]{}
Em um capacitor de $C=\SI{100}{\micro\farad}$, a corrente em $0^+$ é
$\SI{20}{\milli\ampere}$ entrando no terminal positivo. Determine
$\mathrm{d}v_C/\mathrm{d}t|_{0^+}$.
\resposta{$\mathrm{d}v_C/\mathrm{d}t|_{0^+}=\SI{200}{\volt\per\second}$}
\resolucao{Pela convenção passiva, $i_C=C\dot v_C$. Logo
$\dot v_C=0{,}02/(\pot{100}{-6})=\SI{200}{\volt\per\second}$.}
\end{questao}

\begin{questao}[2]{}
Um circuito RLC em paralelo possui $R=\SI{200}{\ohm}$, $C=\SI{100}{\micro\farad}$ e
$L=\SI{0.1}{\henry}$. Classifique a resposta natural da tensão.
\resposta{Subamortecida}
\resolucao{No circuito RLC em paralelo, $\alpha=1/(2RC)=1/[2\cdot200\cdot\pot{100}{-6}]
=\SI{25}{\per\second}$. Já $\omega_0\approx\SI{316.2}{\radian\per\second}$.
Como $\alpha<\omega_0$, a resposta é subamortecida.}
\end{questao}

\begin{questao}[2]{}
Em uma resposta de primeira ordem, qual erro percentual em relação ao valor final
permanece após $5\tau$?
\resposta{$\approx\SI{0.674}{\percent}$}
\resolucao{O erro normalizado é $e^{-t/\tau}$. Em $5\tau$,
$e^{-5}\approx0{,}00674$, ou \SI{0.674}{\percent}.}
\end{questao}

\begin{questao}[2]{}
Um circuito LC ideal possui um capacitor de $C=\SI{100}{\micro\farad}$
com tensão inicial de $\SI{10}{\volt}$ e um indutor de
$L=\SI{50}{\milli\henry}$ com corrente inicial nula.
Determine a energia inicial e a corrente máxima ideal.
\resposta{$E=\SI{5}{\milli\joule}$; $I_{\max}\approx\SI{0.447}{\ampere}$}
\resolucao{$E=\tfrac12CV^2=\tfrac12\pot{100}{-6}\cdot100=\SI{5}{\milli\joule}$.
Quando toda a energia está no indutor,
$\tfrac12LI_{\max}^2=E$, então
$I_{\max}=V\sqrt{C/L}=10\sqrt{\pot{100}{-6}/0{,}05}\approx\SI{0.447}{\ampere}$.}
\end{questao}

\begin{questao}[2]{}
Para $L=\SI{10}{\milli\henry}$ e $C=\SI{100}{\micro\farad}$, determine o
período natural ideal $T_0$.
\resposta{$T_0\approx\SI{6.28}{\milli\second}$}
\resolucao{$\omega_0=1/\sqrt{LC}=1/\sqrt{\pot{1}{-6}}=\SI{1000}{\radian\per\second}$.
Logo $T_0=2\pi/\omega_0\approx\SI{6.28}{\milli\second}$.}
\end{questao}

\gabarito[Nível 2]
